\section{Introduction}
\label{sec:intro}

Terrestrial life evolved over nearly four billion years within the omnipresent, protective envelope of Earth's geomagnetic field (GMF), which delivers a continuous surface magnetic flux density ranging from $\sim\SI{25}{\micro\tesla}$ at the equator to $\sim\SI{65}{\micro\tesla}$ at the magnetic poles. On a planetary scale, the GMF serves as a primary shield against the solar wind and galactic cosmic radiation (GCR). On a cellular and physiological scale, the GMF acts as a pervasive environmental signal that has shaped fundamental biological mechanisms across all domains of life \citep{Maffei2014, Binhi2017}. From cryptochrome-mediated light sensing and radical pair spin recombination kinetics \citep{Hore2016, Ritz2000} to polar auxin transport in plant root gravitropism \citep{Narayana2018}, circadian rhythm entrainment \citep{Yoshii2009}, and cellular redox homeostasis \citep{Belyavskaya2004}, biological systems are intrinsically calibrated to Earth's magnetic baseline.

As humanity transitions from low Earth orbit (LEO) to sustained multi-planetary exploration, space missions will permanently detach biological systems from this terrestrial magnetic cradle. In deep space, on the lunar surface, and across Martian landscapes, organisms will experience chronic exposure to hypomagnetic fields (HMF)—defined as any magnetic environment with an intensity significantly below Earth's baseline ($<\SI{5}{\micro\tesla}$). While interplanetary transit exposes payloads to isotropic interplanetary magnetic fields (IMF) of $\sim\SI{5}{\nano\tesla}$, surface operations on the Moon and Mars present complex, highly heterogeneous planetary magnetic environments resulting from the ancient extinction of their internal core dynamos \citep{Acuna1999, Wieczorek2006}.

The paleomagnetic histories of the Moon and Mars diverged dramatically from Earth's continuously active core dynamo:
\begin{enumerate}
    \item \textbf{The Moon:} The early lunar core dynamo ceased operation between $\sim 3.5$ and $1.9$ billion years ago \citep{Wieczorek2006}. Today, the Moon lacks any global dipole field, preserving only localized, weak crustal magnetic anomalies concentrated primarily in the far-side South Pole-Aitken (SPA) basin and impact basin antipodes \citep{Hood2020, Tsunakawa2015}.
    \item \textbf{Mars:} Mars experienced a catastrophic core dynamo shutdown approximately $4.1$ to $3.9$ billion years ago during the late Noachian era \citep{Acuna1999, Connerney2005}. Unlike the Moon, Mars retains extremely intense, large-scale remnant crustal magnetization embedded within its ancient southern highland crust \citep{Langlais2019}. This creates a profound planetary dichotomy between demagnetized northern lowlands ($|\mathbf{B}_{\text{surf}}| < \SI{25}{\nano\tesla}$) and intensely magnetized southern anomaly belts ($|\mathbf{B}_{\text{surf}}| > \SI{1500}{\nano\tesla}$) capable of generating mini-magnetospheres \citep{Mitchell2001, Johnson2020}.
\end{enumerate}

Despite extensive space biology research into microgravity and space radiation, the physiological consequences of planetary magnetic field heterogeneity have remained largely unexamined in integrated mission planning. Controlled terrestrial HMF simulations reveal that magnetic deprivation induces significant molecular and phenotypic stress, including delayed flowering and altered gravitropism in model crops (\textit{Arabidopsis thaliana}) \citep{Agliassa2018, Narayana2018}, dysregulation of iron uptake \citep{Narayana2021}, restructuring of gut bacterial communities \citep{Zhang2017}, and accelerated trabecular bone loss in mammalian models \citep{Xu2012, Mo2014}.

Here, we present a unified, comparative planetary magnetobiology framework for deep-space exploration. We extract and model the crustal magnetic environment across 36 prospective and historical landing sites—encompassing 23 lunar locations (including all 13 candidate Artemis III South Polar regions, Apollo sites, Chang'e, and Chandrayaan-3) and 13 Martian locations (including ice-rich human ISRU landing sites, active rovers, InSight ground truth, and southern anomaly belts). We synthesize these physical determinations with a systematic review of terrestrial magnetobiology literature to evaluate multi-taxonomic vulnerabilities under the triple-stressor triad (hypomagnetism, altered gravity, and cosmic radiation), establishing actionable engineering and siting recommendations for future planetary habitats.
